\documentclass[11pt,a4paper]{article}
\usepackage[T1]{fontenc}
\usepackage[utf8]{inputenc}
\usepackage{lmodern}
\usepackage{amsmath,amssymb,amsthm,mathtools}
\usepackage[margin=27mm,headheight=14pt]{geometry}
\usepackage{microtype}
\usepackage{enumitem}
\usepackage{xcolor}
\usepackage{fancyhdr}
\usepackage{booktabs}
\usepackage{longtable}
\usepackage{xurl}
\usepackage{hyperref}
\usepackage{bookmark}
\usepackage[capitalise,nameinlink,noabbrev]{cleveref}
\definecolor{linkblue}{RGB}{22,68,105}
\hypersetup{colorlinks=true,linkcolor=linkblue,citecolor=linkblue,urlcolor=linkblue,
  pdftitle={Entropy concavity for four, five and six positively weighted Bernoulli summands},
  pdfauthor={Yitzchak Shmalo},
  pdfsubject={Shannon entropy, weighted Bernoulli sums, exact certificates}}
\urlstyle{same}
\setlength{\parindent}{1.2em}
\setlength{\parskip}{0.2em}
\setlength{\emergencystretch}{1.5em}
\setlist[enumerate]{itemsep=0.35em,topsep=0.45em}
\numberwithin{equation}{section}
\newtheorem{theorem}{Theorem}[section]
\newtheorem{lemma}[theorem]{Lemma}
\newtheorem{proposition}[theorem]{Proposition}
\newtheorem{corollary}[theorem]{Corollary}
\newtheorem{conjecture}[theorem]{Conjecture}
\theoremstyle{remark}
\newtheorem{remark}[theorem]{Remark}
\newcommand{\Pp}{\mathbb P}
\newcommand{\eps}{\varepsilon}
\newcommand{\R}{\mathbb R}
\newcommand{\Z}{\mathbb Z}
\newcommand{\sgn}{\operatorname{sgn}}
\pagestyle{fancy}
\fancyhf{}
\fancyhead[L]{\small Weighted Bernoulli entropy}
\fancyhead[R]{\small 25 November 2026}
\fancyfoot[C]{\thepage}
\renewcommand{\headrulewidth}{0.4pt}

\title{Entropy concavity for four, five and six positively weighted Bernoulli summands}
\author{Yitzchak Shmalo\\[0.4em]
\small Einstein Institute of Mathematics, The Hebrew University of Jerusalem\\
\small Givat Ram, Jerusalem, Israel}
\date{25 November 2026}

\begin{document}
\maketitle
\thispagestyle{fancy}

\begin{abstract}
J\"org Neunh\"auserer conjectures that the entropy of a recursively overlapped Bernoulli distribution is concave in its parameter. Amir asks whether the entropy of $\sum_{i\le m}w_iB_i$, with positive real weights $w_i$ and independent Bernoulli$(p)$ variables $B_i$, is concave in $p$, and proves that it is for $m\le3$. We prove strict concavity for $m=4,5,6$; with Amir's result this gives Neunh\"auserer's conjecture for the first six levels of the recursion. The entropy depends on the weights only through which differences of two subset sums vanish, and this leaves $23$, $216$ and $6{,}767$ cases for $m=4,5,6$; for $m=6$ they come from $23{,}273$ spans of vanishing relations, up to permutation of the weights. For each case an interval $[\delta_m,1/2]$, with $2^{-15}\le\delta_m\le2^{-12}$, is covered by finitely many intervals with rational endpoints, on each of which the second derivative has a negative rational upper bound. On $(0,\delta_m]$ the second derivative is negative by an inequality valid for every $m\ge3$. For every $m$, $p$ times the second derivative tends to $-m$ as $p\downarrow0$.
\end{abstract}

\noindent\textit{Keywords:} Shannon entropy, weighted Bernoulli sums, entropy concavity, subset sums, rational certificates.\\
\textit{2020 Mathematics Subject Classification:} 60C05, 94A17, 65G20.

\section{Introduction}\label{sec:intro}

Let $m$ be a positive integer, let $w=(w_1,\ldots,w_m)$ have positive real coordinates, and let $B_1,\ldots,B_m$ be independent Bernoulli variables with the same parameter $p\in[0,1]$. Set
\begin{equation}\label{eq:sum}
 S_w(p)=\sum_{i=1}^m w_i B_i,\qquad F_w(p)=H(S_w(p)),
\end{equation}
where $H$ is Shannon entropy with natural logarithms~\cite{Shannon}.
The positive-weight concavity question asks whether $F_w$ is concave for every $m$ and every such $w$.

J\"org Neunh\"auserer~\cite{Neun} poses an entropy-concavity conjecture for recursively overlapped Bernoulli distributions; the construction interpolates between a binomial law and a Bernoulli product law, and its motivation includes random walks on self-similar graphs. Let $(\mathfrak a_n)$ be a strictly increasing sequence of integers with $\mathfrak a_1=2$, put $P_1(p)=(p,1-p)$ and $P_{n+1}(p)=p\mathcal L(P_n(p))+(1-p)\mathcal R(P_n(p))$, where $\mathcal L$ and $\mathcal R$ place a vector of length $\mathfrak a_n$ at the left and at the right end of a vector of length $\mathfrak a_{n+1}$. Neunh\"auserer~\cite[Conjecture~3.1]{Neun} conjectures that the entropy of $P_n(p)$ is concave in $p\in[0,1]$ for every such sequence and every $n$. The law $P_n(p)$ is that of $\sum_{i=1}^n \omega_iB_i'$ with $\omega_1=1$, $\omega_i=\mathfrak a_i-\mathfrak a_{i-1}$ and independent Bernoulli$(1-p)$ variables $B_i'$. Since $p\mapsto1-p$ preserves concavity, the conjecture for a given $n$ is the case $m=n$ of the positive-weight concavity question for integer weights one of which equals $1$. Amir~\cite{Amir} gives the positive-weight formulation~\eqref{eq:sum}, with the shifts $\omega_i$ as weights; integer weights include the recursive construction, and arbitrary positive real weights extend it. Neunh\"auserer~\cite[Proposition~3.1]{Neun} proves the conjecture for $\mathfrak a_n=n+1$, where $P_n(p)$ is a shifted binomial law, and for sequences with $\mathfrak a_{n+1}\ge2\mathfrak a_n$, where the $2^n$ subset sums are distinct and the entropy is $n$ times the binary entropy.

Amir~\cite{Amir} proves that $F_w$ is strictly concave for $m\le3$; the case $m=1$ is the binary entropy function. For $m=4$ he sorts the weight vectors by which subset sums coincide, which up to permutation of the weights gives $23$ classes. He proves seven of them analytically and checks all $23$ numerically; for $m=4$, \cref{thm:main} below turns this check into a proof for the remaining sixteen classes.

Three things stand in the way of extending this case analysis; the first two arise from $m=5$ on, the third already for $m=4$. First, the weights are real, and rounding them can separate two subset sums that are equal, which changes the law of $S_w(p)$. Second, the number of cases grows. For a value $s$ of $S_w(p)$ let $N_{s,k}$ be the number of $k$-element sets $A\subseteq\{1,\ldots,m\}$ with $\sum_{i\in A}w_i=s$. Then $S_w(p)=s$ has probability $\sum_kN_{s,k}p^k(1-p)^{m-k}$, so $F_w$ depends on $w$ only through the multiset of the vectors $(N_{s,0},\ldots,N_{s,m})$, the \emph{count vectors} of the values $s$; this multiset is the \emph{entropy profile} of $w$. The values of $S_w(p)$ are the subset sums of $w$. There are $23$ entropy profiles for $m=4$, one for each of Amir's classes, $216$ for $m=5$ and $6{,}767$ for $m=6$ (\cref{tab:profiles}, \cref{prop:sixprofiles}). Third, an upper bound for $F_w''$ on an interval of values of $p$, obtained by bounding each term, needs a positive lower bound on that interval for the probability of every value $s$. For $s\ne0$ this probability vanishes at $p=0$, and $F_w''$ is unbounded there: $pF_w''(p)\to-m$ as $p\downarrow0$ (\cref{thm:limit}).

For the first obstruction, every positive $w$ has an integer vector $v$ with
\[
 1\le v_i\le\lfloor m^{m/2}\rfloor
 \quad\text{and}\quad
 \sgn\Bigl(\sum_{i\in A}v_i-\sum_{i\in C}v_i\Bigr)
 =\sgn\Bigl(\sum_{i\in A}w_i-\sum_{i\in C}w_i\Bigr)
 \;\text{for all }A,C\subseteq\{1,\ldots,m\}
\]
(\cref{lem:representative}). The vector $v$ is an integer multiple of a vertex of a polyhedron defined by these signs, and Cramer's rule and Hadamard's inequality bound its coordinates. Saburo Muroga, Iwao Toda and Satoru Takasu~\cite{MTT} prove that every linear threshold function of $n$ Boolean variables, that is, every function $x\mapsto\sgn(\sum_iu_ix_i-t)$ on $\{0,1\}^n$ that is never zero, has a representation with integer weights of absolute value $2^{O(n\log n)}$ (see the account of Rocco Servedio~\cite{Servedio}). The usual proof of bounds of this kind takes the weights at a vertex of the polyhedron of admissible weights and bounds them by Cramer's rule and Hadamard's inequality; \cref{lem:representative} applies this argument to the system of all comparisons of two subset sums of $w$, in which the sign $0$ occurs, and keeps the comparisons that are equalities.

For the second obstruction, the sign $0$ is kept, so $v$ and $w$ have the same entropy profile, and for $m=4$ and $m=5$ the entropy profiles are found by listing the integer vectors with coordinates at most $16$ and $55$. For $m=6$ the bound is $216$ and the integer vectors are too many to list. There the entropy profile of $w$ is determined by the linear span of the vectors $a\in\{-1,0,1\}^6$ with $a\cdot w=0$, the \emph{relation space} of $w$, and up to permutation of the coordinates there are $23{,}273$ relation spaces (\cref{lem:relations}, \cref{tab:six}).

For the third obstruction, the interval $(0,1/2]$ is cut at the point $\delta_m$, where $\delta_4=2^{-12}$, $\delta_5=2^{-14}$ and $\delta_6=2^{-15}$. For each entropy profile the interval $[\delta_m,1/2]$ is divided into finitely many intervals with rational endpoints, and on each of them $F_w''$ is bounded above by a negative rational number computed in rational arithmetic (\cref{prop:certificates}). On $(0,\delta_m]$ the sign comes from one value of $S_w(p)$. A set $A$ with $\sum_{i\in A}w_i=\min_iw_i$ has one element, so $S_w(p)=\min_iw_i$ has probability $rp(1-p)^{m-1}$, where $r$ is the number of weights equal to the smallest. From this formula, and since every value $s$ has probability at least $p^m$ when $p\le1/2$, one obtains, for every $m\ge3$,
\[
 pF_w''(p)\le 4m^2(m-1)\,p\log\frac1p-(1-p)^{m-3}(1-mp)^2\qquad\Bigl(0<p\le\frac12\Bigr)
\]
(\cref{lem:endpoint}). Hence $F_w''(p)<-1/(2p)$ for $0<p\le1/(64m^4)$ (\cref{thm:allendpoint}), and $F_w''<0$ on $(0,\delta_m]$ for $m=4,5,6$ (\cref{cor:endpoint}). The symmetry $F_w(p)=F_w(1-p)$ gives the interval $[1/2,1)$. With Amir's theorem for $m\le3$, \cref{thm:main} gives Neunh\"auserer's conjecture for $n\le6$. \Cref{thm:main} is new for $m=5,6$ and, for $m=4$, for the sixteen classes that Amir checks only numerically; the question remains open for $m\ge7$, where \cref{thm:allendpoint,thm:limit} still hold.

\begin{theorem}\label{thm:main}
Let $m\in\{4,5,6\}$ and let $w_1,\ldots,w_m$ be positive real numbers. If $B_1,\ldots,B_m$ are independent Bernoulli$(p)$ random variables, then
\[
 p\longmapsto H\!\left(\sum_{i=1}^m w_i B_i\right)
\]
has strictly negative second derivative on $(0,1)$ and is strictly concave on $[0,1]$.
\end{theorem}

Lawrence Shepp and Ingram Olkin~\cite{SO} formulate the entropy-concavity problem for sums with unit weights. Erwan Hillion and Oliver Johnson~\cite{HJ} prove joint concavity in the individual Bernoulli parameters in that setting. Their theorem does not apply to integer weights: $w_iB_i$ with an integer $w_i>1$ is not a sum of $w_i$ independent unit-weight Bernoulli variables. Jiange Li~\cite{Li} compares R\'enyi entropies of different orders for weighted Bernoulli sums. Haoran Wang~\cite{Wang} studies R\'enyi and Tsallis concavity for unit-weight sums. None of these results gives concavity of the Shannon entropy of a weighted sum in a common parameter.

Hadamard's inequality bounds the determinants in \cref{lem:representative}~\cite{Matrix}. For six summands, divisibility of determinants of matrices with entries in $\{-1,1\}$ by $2^{m-1}$ lowers the coordinate bound to $192$ (\cref{cor:refinement}), but the nondecreasing vectors in $\{1,\ldots,192\}^6$ still number $\binom{197}{6}$, and the enumeration of relation spaces replaces the listing of these integer vectors. The hyperplanes $a\cdot w=0$ with $a\in\{-1,0,1\}^m$ form the arrangement of comparative probability theory: Ilya Chevyrev, Dominic Searles and Arkadii Slinko~\cite[Section~2.4]{CSS} describe it in the positive simplex, where weight vectors normalized to sum one lie. An entropy profile depends only on which of these hyperplanes contain $w$, so the relation spaces enumerated here are the orthogonal complements of the flats of that arrangement that meet the open positive orthant. The interval bounds use Sergei Bernstein's polynomial basis and Paul de Casteljau's subdivision algorithm, as presented by Rida Farouki~\cite{Farouki}, together with rational bounds for the logarithm. When the weights take values $\theta_1,\ldots,\theta_\ell$ with multiplicities $r_1,\ldots,r_\ell$ and the weighted sum determines how many Bernoulli variables of each weight succeed, the entropy is the sum of the entropies of $\operatorname{Bin}(r_j,p)$, which are concave~\cite{SO,HJ}. \Cref{thm:blocks} records this consequence, with an elementary proof of a strict bound, and gives strict concavity for every number of summands. The case $\ell=1$ is the binomial law and the case $r_1=\cdots=r_\ell=1$ is that of $2^m$ distinct subset sums, which contain the two cases of~\cite[Proposition~3.1]{Neun}.

\section{Setup and notation}\label{sec:setup}

Write $[m]=\{1,\ldots,m\}$ and let $\mathcal S_w$ be the finite set of subset sums of $w$. With the counts $N_{s,k}$ introduced in \cref{sec:intro}, put
\begin{equation}\label{eq:counts}
 q_s(p)=\sum_{k=0}^m N_{s,k}p^k(1-p)^{m-k}\qquad(s\in\mathcal S_w).
\end{equation}
The fibre of $s$ is the set of $A\subseteq[m]$ with $\sum_{i\in A}w_i=s$. Writing $\Pp$ for probability under the law of the independent Bernoulli variables, $q_s(p)=\Pp(S_w(p)=s)$ and
\begin{equation}\label{eq:entropy}
 F_w(p)=-\sum_{s\in\mathcal S_w}q_s(p)\log q_s(p),
\end{equation}
with $0\log0=0$. In the entropy profile defined in \cref{sec:intro}, equal count vectors occur with their multiplicities. The multiset determines $F_w$ for every $p$, even though it does not retain the numerical values of the subset sums. Weight vectors with different signs of differences of two subset sums can have the same entropy profile.

For $0<p<1$ all $q_s$ are positive, so differentiation gives
\begin{equation}\label{eq:curvature}
 F_w''(p)=-\sum_s q_s''(p)\log q_s(p)
            -\sum_s\frac{(q_s'(p))^2}{q_s(p)}.
\end{equation}
Here $\sum_s q_s''=0$ removes the term $-\sum_s q_s''$ from the second derivative. Complementing every subset gives
\begin{equation}\label{eq:symmetry}
 q_s(p)=q_{W-s}(1-p),\qquad F_w(p)=F_w(1-p),
 \qquad W=\sum_{i=1}^m w_i.
\end{equation}
The function $F_w$ is continuous on $[0,1]$, with $F_w(0)=F_w(1)=0$.

\section{Integer vectors with the same signs of differences of subset sums}\label{sec:finite}

\begin{lemma}\label{lem:representative}
For every integer $m\ge1$ and every $w\in\R^m$ with positive coordinates, there exists $v\in\Z^m$ such that
\[
 1\le v_i\le M_m:=\lfloor m^{m/2}\rfloor\quad(1\le i\le m)
\]
and, for every two subsets $A,C\subseteq[m]$,
\[
 \sgn\!\left(\sum_{i\in A}v_i-\sum_{i\in C}v_i\right)
 =\sgn\!\left(\sum_{i\in A}w_i-\sum_{i\in C}w_i\right).
\]
In particular, $v$ and $w$ have the same entropy profile.
\end{lemma}

\begin{proof}
Every difference of subset indicators is a vector in $\{-1,0,1\}^m$, and every vector in that set is such a difference. For each nonzero vector $a$ with $a\cdot w=0$, impose $a\cdot x=0$. For each other nonzero $a$, choose its sign so that $a\cdot w>0$ and impose $a\cdot x\ge1$. Also impose $x_i\ge1$ for every $i$. Let $\mathcal Q$ be the resulting polyhedron. There are finitely many constraints, and a sufficiently large positive multiple of $w$ belongs to $\mathcal Q$, so $\mathcal Q$ is nonempty.

Minimize $\sum_i x_i$ over $\mathcal Q$. A nonempty sublevel set is closed and bounded because $x_i\ge1$, hence compact. The minimum is attained, and its minimizing face is nonempty and compact. Choose an extreme point $x^*$ of this face. If $x^*$ were a nontrivial convex combination of two points of $\mathcal Q$, minimality would force both points into the minimizing face. Thus $x^*$ is a vertex of $\mathcal Q$.

The normals of the equalities and the inequalities tight at $x^*$ span $\R^m$. Otherwise a nonzero vector orthogonal to all those normals would permit sufficiently small perturbations of $x^*$ in both directions, preserving the equalities, tight inequalities and all the finitely many slack inequalities. This would contradict extremality. Select $m$ independent tight constraints, and write them as $Gx^*=b$. The entries of $G$ lie in $\{-1,0,1\}$, the entries of $b$ lie in $\{0,1\}$, and $D:=\det G$ is a nonzero integer.

Let $G_i$ be obtained from $G$ by replacing column $i$ with $b$. Cramer's rule gives
\[
 |D|x_i^*=\sgn(D)\det G_i.
\]
Every entry of $G_i$ still has absolute value at most $1$. Hadamard's inequality therefore gives $|\det G_i|\le m^{m/2}$. Consequently $v:=|D|x^*$ has positive integer coordinates bounded by $M_m$.
Every imposed equality remains zero after multiplication by $|D|$, and every imposed inequality has positive left side. Every difference of two subset sums has the form $a\cdot x$ with $a\in\{-1,0,1\}^m$, so the signs agree. The equality relations identify exactly the same subsets in each fibre, preserving their cardinalities and hence the entropy profile.
\end{proof}

\begin{corollary}\label{cor:refinement}
Every positive real weight vector for six summands has a positive integer vector with coordinates at most $192$ preserving the signs of all differences of two subset sums.
\end{corollary}
\begin{proof}
Use the vertex and Cramer numerators in \cref{lem:representative}. Multilinearity, applied one entry at a time, bounds the absolute determinant of every matrix with entries in $[-1,1]$ by that of some matrix with entries in $\{-1,1\}$. For a matrix with entries in $\{-1,1\}$ of order $m$, subtract its first row from the other $m-1$ rows. Each of those rows has even entries, so its determinant is divisible by $2^{m-1}$. At order six Hadamard's inequality gives $216$, and hence the determinant is at most
\[
 32\lfloor216/32\rfloor=192.
\]
This bounds every Cramer numerator and proves the assertion.
\end{proof}

The bound $192$ uses only divisibility of determinants of matrices with entries in $\{-1,1\}$. The stronger determinant estimate of Ehlich and Wojtas~\cite{Ehlich,Wojtas} gives $160$ at order six and can be substituted in the same Cramer argument.

Permutation of coordinates preserves the entropy profile. Therefore it suffices to enumerate nondecreasing integer vectors with coordinates between $1$ and $M_m$. Rational approximation alone would not justify this reduction: it can separate subsets whose sums are exactly equal.

\subsection{Enumeration of entropy profiles}\label{sec:enumeration}

For a fixed integer vector, the integer sums partition its $2^m$ subsets into fibres. The count vectors of these fibres, sorted lexicographically, determine its entropy profile. Listing every nondecreasing integer vector with coordinates in $\{1,\ldots,M_m\}$ gives \cref{tab:profiles}.

\begin{table}[ht]
\centering
\begin{tabular}{rrrr}
\toprule
$m$ & $M_m$ & integer vectors & entropy profiles\\
\midrule
$4$ & $16$ & $3{,}876$ & $23$\\
$5$ & $55$ & $5{,}006{,}386$ & $216$\\
\bottomrule
\end{tabular}
\caption{Listing of the integer vectors with coordinates at most $M_m$. The vector count is $\binom{M_m+m-1}{m}$.}\label{tab:profiles}
\end{table}

\Cref{tab:four} lists one positive integer vector for each four-summand entropy profile. Each vector has the smallest coordinate sum among the listed integer vectors with its profile, and comes first in lexicographic order among those.

\begin{table}[ht]
\centering\small
\begin{tabular}{crcrcr}
\toprule
weights & $|\mathcal S_v|$ & weights & $|\mathcal S_v|$ & weights & $|\mathcal S_v|$\\
\midrule
$(1,1,1,1)$ & $5$ & $(1,1,3,3)$ & $9$ & $(1,3,3,5)$ & $11$\\
$(1,1,1,2)$ & $6$ & $(1,2,2,3)$ & $9$ & $(1,1,3,6)$ & $12$\\
$(1,1,1,3)$ & $7$ & $(1,1,2,5)$ & $10$ & $(1,2,3,5)$ & $12$\\
$(1,1,2,2)$ & $7$ & $(1,1,3,4)$ & $10$ & $(1,2,3,6)$ & $13$\\
$(1,1,1,4)$ & $8$ & $(1,2,3,3)$ & $10$ & $(1,2,4,5)$ & $13$\\
$(1,1,2,3)$ & $8$ & $(1,1,3,5)$ & $11$ & $(1,2,3,7)$ & $14$\\
$(1,1,2,4)$ & $9$ & $(1,2,3,4)$ & $11$ & $(1,2,4,7)$ & $15$\\
 & & & & $(1,3,5,7)$ & $15$\\
 & & & & $(1,2,4,8)$ & $16$\\
\bottomrule
\end{tabular}
\caption{The four-summand entropy profiles, with positive integer weights $v$ and the number $|\mathcal S_v|$ of subset sums.}\label{tab:four}
\end{table}

\section{Six summands: relation-space orbits}\label{sec:six}

Let $T_m$ contain one of each pair $a,-a$ of mixed-sign vectors in $\{-1,0,1\}^m$; its elements are called ternary normals. For $w>0$ define its \emph{relation space}
\[
 L(w)=\operatorname{span}\{a\in T_m:a\cdot w=0\}.
\]
A subspace spanned by ternary normals is called feasible if its orthogonal complement contains a vector with all coordinates strictly positive. Its rank is its vector-space dimension. The orbit of a subspace is the set of its images under permutations of the coordinates. For six summands, $T_6$ has
$(3^6-2\cdot2^6+1)/2=301$ elements.

\begin{lemma}\label{lem:relations}
For every positive real weight vector $w$ for $m$ summands, the feasible relation space $L(w)$ has rank at most $m-1$, and two subsets $A,C\subseteq[m]$ satisfy $\sum_{i\in A}w_i=\sum_{i\in C}w_i$ if and only if $\mathbf 1_A-\mathbf 1_C\in L(w)$. Conversely, every feasible space spanned by ternary normals is the relation space of a positive real vector. Starting with the zero space and adjoining each ternary normal to every feasible space, discarding infeasible extensions, reaches every such space. Retaining one subspace from each orbit under coordinate permutations reaches every orbit.
\end{lemma}
\begin{proof}
For subsets $A,C\subseteq[m]$ with indicator vectors $\mathbf 1_A,\mathbf 1_C$, the equality $\sum_{i\in A}w_i=\sum_{i\in C}w_i$ is equivalent to $\mathbf 1_A-\mathbf 1_C\in L(w)$. A nonzero difference that vanishes on a positive vector has both signs, so one orientation is a ternary normal; the reverse implication follows from $L(w)\subseteq w^\perp$.

For a feasible space $L$ spanned by ternary normals, the positive part of $L^\perp$ is nonempty and relatively open. For $a\in T_m\setminus L$, the hyperplane $a^\perp\cap L^\perp$ is proper. Avoiding these finitely many hyperplanes gives a positive vector $w$ with $L(w)=L$. A chain of ternary normals forming a basis of $L$ reaches $L$; each intermediate orthogonal complement contains the same positive vector. Finally, coordinate permutations preserve the set of ternary normals, feasibility and subset cardinalities. Extending one subspace from each orbit by every ternary normal therefore reaches every successor orbit.
\end{proof}

The maximal minors $P_{\mathcal I}$ of an integer row basis of a rank-${\varrho}$ subspace, indexed by the ${\varrho}$-element coordinate subsets $\mathcal I$, determine the subspace up to a common nonzero factor. Dividing by their common divisor and making the first nonzero coordinate positive gives its primitive exterior coordinates; at rank zero the single coordinate is $1$. Adjoining a ternary normal takes an exterior product, which is zero exactly when the normal already belongs to the subspace. Coordinate permutations act by signed permutations of the minors, so the lexicographically least oriented vector determines one subspace from each orbit. The maximum absolute primitive coordinates at ranks zero through five are $1,1,2,4,12,30$.

\begin{lemma}\label{lem:positive}
Let $L\subseteq\R^m$ have rank ${\varrho}<m$, with maximal minors $P_{\mathcal I}$ of a row basis. For each ordered coordinate set $J=\{j_0<\cdots<j_{\varrho}\}$, form the vector supported on $J$ whose entries are
\[
 \kappa_{j_\iota}=(-1)^{\iota}P_{J\setminus\{j_\iota\}}\quad(0\le\iota\le\varrho).
\]
Retain each nonzero vector whose entries have one sign, oriented nonnegatively. Then $L^\perp$ contains a strictly positive vector if and only if the retained supports cover all coordinates. If it does, the retained vectors span $L^\perp$.
\end{lemma}
\begin{proof}
Cofactor expansion puts every listed vector in $L^\perp$. Let $\mathcal K=L^\perp\cap[0,\infty)^m$. It is a pointed polyhedral cone. Intersecting with $\sum_i x_i=1$ shows that it is generated by its extreme rays. For an extreme ray with support $E$, the corresponding columns of a row basis have rank $|E|-1$: lower rank would permit a nonproportional supported perturbation in both directions, and higher rank would leave no nonzero supported kernel vector. Extend these columns to rank ${\varrho}$ using ${\varrho}-|E|+1$ further columns. On the resulting ${\varrho}+1$ columns, the kernel is one-dimensional. The cofactor vector is nonzero, agrees with the ray on $E$ up to scale, and is zero on the added columns. Thus the retained list contains every extreme ray.

A strictly positive vector in $\mathcal K$ implies that every coordinate occurs in some extreme ray. Conversely, summing the retained nonnegative vectors gives a strictly positive vector when their supports cover. If $x>0$ belongs to $\mathcal K$, then $x+\eps\mathbf y\in \mathcal K$ for every $\mathbf y\in L^\perp$ and sufficiently small $\eps>0$. Hence $\mathcal K$ spans $L^\perp$, and so do the retained vectors.
\end{proof}

For a positive weight vector $w$ with $L(w)=L$, two subsets have equal sums precisely when their indicator vectors have equal scalar products with every retained vector. The fibres are therefore determined by these scalar products, and their count vectors give the entropy profile. The enumeration of relation spaces gives \cref{tab:six}. For $m=4$ and $m=5$ the same enumeration gives $23$ and $297$ relation spaces up to permutation of coordinates and the $23$ and $216$ entropy profiles of \cref{tab:profiles}.

\begin{table}[ht]
\centering
\begin{tabular}{rrr}
\toprule
rank & relation-space orbits & entropy profiles through this rank\\
\midrule
$0$ & $1$ & $1$\\
$1$ & $9$ & $10$\\
$2$ & $117$ & $76$\\
$3$ & $1{,}593$ & $556$\\
$4$ & $10{,}334$ & $2{,}948$\\
$5$ & $11{,}219$ & $6{,}767$\\
\bottomrule
\end{tabular}
\caption{The enumeration of relation spaces for six summands: $23{,}273$ relation-space orbits and $6{,}767$ entropy profiles.}\label{tab:six}
\end{table}

\begin{proposition}\label{prop:sixprofiles}
The sums of six independent Bernoulli$(p)$ variables with positive real weights have exactly $6{,}767$ entropy profiles. Each profile is the entropy profile of a positive integer weight vector with coordinates at most $36$.
\end{proposition}
\begin{proof}
By \cref{lem:relations,lem:positive}, the enumeration of relation spaces, which starts at the zero space and adjoins ternary normals, lists every relation space $L(w)$, $w\in(0,\infty)^6$, up to permutation of coordinates. There are $23{,}273$ of them, by rank as in \cref{tab:six}, and they have $6{,}767$ distinct entropy profiles. One positive integer vector with coordinates at most $36$ for each of these profiles is in the file \texttt{profiles\_six\_compact.json} of the repository (Code and data availability). Such a vector need not preserve the signs of all differences of two subset sums of another vector with that profile.
\end{proof}


\section{Uniform endpoint bounds}\label{sec:endpoint}

\begin{lemma}\label{lem:endpoint}
Let $m\ge3$ and let $w\in\R^m$ have positive coordinates. For the entropy function in~\eqref{eq:sum}, one has, for $0<p\le1/2$,
\begin{equation}\label{eq:endpoint}
 pF_w''(p)\le4m^2(m-1)p\log(1/p)
              -(1-p)^{m-3}(1-mp)^2.
\end{equation}
\end{lemma}

\begin{proof}
The probability of every subset of successful Bernoulli variables is at least $p^m$ when $p\le1/2$. Every nonempty fibre contains such a subset, so $p^m\le q_s(p)\le1$ and $|\log q_s(p)|\le m\log(1/p)$.

Write the product law of the Bernoulli variables as
\[
 \mu_p=\bigotimes_{i=1}^m\bigl((1-p)\mathbf e_0+p\mathbf e_1\bigr),
\]
where $\mathbf e_j$ is the unit point mass at $j$. The second derivative of $\mu_p$ is a sum over $m(m-1)$ ordered pairs of distinct coordinates. Each summand has two factors $\mathbf e_1-\mathbf e_0$, each of $\ell^1$ norm $2$, and probability measures in the other factors. Its $\ell^1$ norm is $4$. Summing a signed measure over the fibres of a map cannot increase its $\ell^1$ norm. Therefore
\begin{equation}\label{eq:norm}
 \sum_s|q_s''(p)|\le4m(m-1),\qquad
 \left|\sum_s q_s''(p)\log q_s(p)\right|
 \le4m^2(m-1)\log(1/p).
\end{equation}

Let $w_{\min}=\min_i w_i$, attained $r\ge1$ times. A subset with at least two elements has sum at least $2w_{\min}>w_{\min}$, so the fibre at $w_{\min}$ consists exactly of these $r$ singleton subsets. Consequently
\[
 q_{w_{\min}}(p)=rp(1-p)^{m-1},\qquad
 \frac{(q_{w_{\min}}'(p))^2}{q_{w_{\min}}(p)}
 =\frac{r(1-p)^{m-3}(1-mp)^2}{p}.
\]
Use this one term in the nonnegative second sum of~\eqref{eq:curvature}, apply~\eqref{eq:norm}, and use $r\ge1$. This proves~\eqref{eq:endpoint}.
\end{proof}

\begin{corollary}\label{cor:endpoint}
For every positive real weight vector for four summands,
$F_w''(p)<-3/(16p)$ on $0<p\le1/4096$. For every positive real weight vector for five summands,
$F_w''(p)<-1/(4p)$ on $0<p\le1/16384$.
For every positive real weight vector for six summands,
$F_w''(p)<-1/(4p)$ on $0<p\le1/32768$.
In each case, the reflected bound holds near $1$ with $p$ replaced by $1-p$.
\end{corollary}

\begin{proof}
The function $p\log(1/p)$ increases on $(0,\delta_m]$, and $\log2<1$. On the same interval the subtracted term $(1-p)^{m-3}(1-mp)^2$ in~\eqref{eq:endpoint} is at least its value at $\delta_m$, because both of its factors are positive and decreasing there. For $m=4$,
\[
192\delta_4\log(1/\delta_4)<\frac9{16},\qquad
(1-\delta_4)(1-4\delta_4)^2
 =\frac{4285536255}{4294967296}>\frac34.
\]
For $m=5$,
\[
400\delta_5\log(1/\delta_5)<\frac{175}{512}<\frac12,\qquad
(1-\delta_5)^2(1-5\delta_5)^2
 =\frac{72004829826842649}{72057594037927936}>\frac34.
\]
For $m=6$,
\[
720\delta_6\log(1/\delta_6)<\frac{675}{2048}<\frac12,\qquad
(1-\delta_6)^3(1-6\delta_6)^2
=\frac{9440410169765069258743}{9444732965739290427392}>\frac34.
\]
Substitution proves the bounds; \eqref{eq:symmetry} gives their reflections.
\end{proof}

\begin{theorem}\label{thm:allendpoint}
Let $m\ge3$, let $w_1,\ldots,w_m$ be positive real numbers, and let $B_i$ be independent Bernoulli$(p)$ variables. Put $\eta_m=1/(64m^4)$. Then
\[
 F_w''(p)<-\frac1{2p}\quad(0<p\le\eta_m),\qquad
 F_w''(p)<-\frac1{2(1-p)}\quad(1-\eta_m\le p<1).
\]
\end{theorem}
\begin{proof}
For $m\ge3$, $\log64<6$ and $\log m<m-1$ give
$\log(1/\eta_m)<4m+2$. Since $p\log(1/p)$ increases below $e^{-1}$, for $0<p\le\eta_m$ the positive term in \cref{lem:endpoint} satisfies
\[
 4m^2(m-1)p\log(1/p)
 <\frac{(m-1)(4m+2)}{16m^2}
 =\frac14-\frac1{8m}-\frac1{8m^2}<\frac14.
\]
Both factors of $(1-p)^{m-3}(1-mp)^2$ are nonnegative and decreasing on this interval. Bernoulli's inequality gives
\[
 (1-\eta_m)^{m-3}(1-m\eta_m)^2
 \ge 1-(3m-3)\eta_m>\frac34.
\]
For the first inequality, use $(1-\eta_m)^{m-3}\ge1-(m-3)\eta_m$ and $(1-m\eta_m)^2\ge1-2m\eta_m$; their lower bounds are nonnegative and their product is at least the displayed expression. For the last inequality, $3m-3<16m^4$. Substitution into \eqref{eq:endpoint} gives $pF_w''<-1/2$. Complementing the Bernoulli variables proves the reflected inequality.
\end{proof}

\begin{theorem}\label{thm:limit}
Let $m\ge1$ and let $w_1,\ldots,w_m$ be positive real numbers. For independent Bernoulli$(p)$ variables, the entropy in~\eqref{eq:sum} satisfies
\[
 \lim_{p\downarrow0}p\frac{d^2}{dp^2}H\!\left(\sum_iw_iB_i\right)=-m,
 \qquad
 \lim_{p\uparrow1}(1-p)\frac{d^2}{dp^2}H\!\left(\sum_iw_iB_i\right)=-m.
\]
\end{theorem}
\begin{proof}
Fix the weights. For each fibre of $s\ne0$ let $k_s\ge1$ be its minimum subset cardinality and let $\zeta_s$ be the number of subsets in that fibre of cardinality $k_s$. The polynomial $q_s$ of that fibre has expansion
$q_s(p)=\zeta_sp^{k_s}+O(p^{k_s+1})$. If $k_s=1$, then
\[
 p\frac{(q_s')^2}{q_s}\to \zeta_s;
\]
if $k_s\ge2$, the same expression is $O(p^{k_s-1})$ and tends to zero. Positivity isolates the fibre of $0$, whose probability is $(1-p)^m$; its contribution to this limit is zero. The coefficients $\zeta_s$ of the fibres with $k_s=1$ total $m$, since every singleton belongs to exactly one such fibre.

For $k_s=1$, $q_s''=O(1)$, so $pq_s''\log q_s=O(p|\log p|)\to0$. For $k_s\ge2$ it is $O(p^{k_s-1}|\log p|)\to0$. For the fibre of $0$, $q_s''$ is bounded and $\log q_s=O(p)$, so its term also vanishes. Applying these limits to the finite sum \eqref{eq:curvature} gives $-m$. Complement symmetry gives the other endpoint.
\end{proof}

\section{Rational certificates on \texorpdfstring{$[\delta_m,1/2]$}{[delta\_m,1/2]}}\label{sec:certificates}

\subsection{Polynomial interval bounds}\label{sec:bernstein}

For a polynomial $f$ of degree at most $d$, write
\begin{equation}\label{eq:bernstein}
 f(p)=\sum_{k=0}^d c_k\binom dk p^k(1-p)^{d-k}.
\end{equation}
On $[0,1]$ the Bernstein basis functions are nonnegative and sum to $1$, so $\min_k c_k\le f\le\max_k c_k$. The coefficients of the derivative in degree $d-1$ are $d(c_{k+1}-c_k)$. In particular, the degree-$m$ coefficients of $q_s$ are $N_{s,k}/\binom mk$.

For rational $\alpha,\beta$ with $0<\alpha<\beta\le1/2$, de Casteljau subdivision~\cite{Farouki} restricts~\eqref{eq:bernstein} to $I=[\alpha,\beta]$ with rational coefficients. One split at $\alpha$ gives the restriction to $[\alpha,1]$, and a second split at $(\beta-\alpha)/(1-\alpha)$ gives the restriction to $I$. At a split parameter $\tau$, the coefficient recursion is
\[
 c_j^{(0)}=c_j,\qquad
 c_j^{(\sigma+1)}=(1-\tau)c_j^{(\sigma)}+\tau c_{j+1}^{(\sigma)}.
\]
The left boundary of this triangular array and the reversed right boundary are the two sets of restricted coefficients. Applying the bounds $\min_kc_k\le f\le\max_kc_k$ to the restrictions of $q_s,q_s',q_s''$ gives rational interval bounds on $I$.

\subsection{Rational logarithm bounds}\label{sec:logs}

For rational $\chi>0$, choose an integer $\nu$ such that $\chi=2^\nu y$ and $1\le y<2$. Put $z=(y-1)/(y+1)$. For an integer $\mathcal N\ge1$, define
\begin{equation}\label{eq:log}
 \mathsf L_{\mathcal N}(y)=2\sum_{j=0}^{\mathcal N-1}\frac{z^{2j+1}}{2j+1},\qquad
 E_{\mathcal N}(y)=\frac{2z^{2\mathcal N+1}}{(2\mathcal N+1)(1-z^2)}.
\end{equation}
Expanding $2\int_0^z(1-u^2)^{-1}\,du$ and bounding the remaining positive series gives
\[
 \mathsf L_{\mathcal N}(y)\le\log y\le \mathsf L_{\mathcal N}(y)+E_{\mathcal N}(y).
\]
Here $0\le z<1/3$. The same formula at $y=2$, where $z=1/3$, encloses $\log2$. Adding $\nu\log2$, with interval endpoints reversed for this product when $\nu<0$, encloses $\log \chi$. All these interval bounds use rational arithmetic.

\subsection{Bounds on \texorpdfstring{$F_w''$}{the second derivative} and the partitions}\label{sec:partitions}

A cell is a closed interval with rational endpoints. Suppose that the polynomial interval bounds on a cell $I$ are
\[
 q_s\in[q_-,q_+],\quad q_s'\in[\dot q_-,\dot q_+],\quad
 q_s''\in[\ddot q_-,\ddot q_+],\qquad q_->0,
\]
and that the rational logarithm bounds of \cref{sec:logs} give rational numbers $\lambda_-\le\lambda_+$ with $\log q_s\in[\lambda_-,\lambda_+]$.
Set
\begin{align*}
 T_-&=\min\{\ddot q_-\lambda_-,\ddot q_-\lambda_+,\ddot q_+\lambda_-,\ddot q_+\lambda_+\},\\
 V_-&=\begin{cases}0,&\dot q_-\le0\le \dot q_+,\\
 \min\{(\dot q_-)^2,(\dot q_+)^2\},&\text{otherwise}.
 \end{cases}
\end{align*}
The contribution of this fibre to~\eqref{eq:curvature} is at most $-T_--V_-/q_+$. Let $\widetilde U_I$ denote the sum of these bounds over all fibres, and let $U_I$ be $\widetilde U_I$ rounded upward to a multiple of $10^{-6}$; for six summands each fibre's bound is first rounded upward to a multiple of $2^{-40}$. Since every rounding is upward, $U_I$ is rational and $F_w''(p)\le\widetilde U_I\le U_I$ for every $p\in I$.

Starting at $[\delta_m,1/2]$, bisect a cell until $U_I<0$. A certificate for an entropy profile is such a partition of $[\delta_m,1/2]$ into rational cells $I$, each with a rational $U_I<0$ bounding $F_w''$ on $I$.

\begin{proposition}\label{prop:certificates}
For each entropy profile of a sum of four independent Bernoulli$(p)$ variables with positive real weights, there is a partition of $[1/4096,1/2]$ into closed rational cells and, for each cell $I$, a rational number $U_I<0$ with $F_w''(p)\le U_I$ for all $p\in I$. The corresponding assertions hold for five summands on $[1/16384,1/2]$ and six summands on $[1/32768,1/2]$. The partitions have the counts and bounds in \cref{tab:certificates}.
\end{proposition}

\begin{table}[ht]
\centering
\begin{tabular}{rrrr}
\toprule
$m$ & profiles & cells in all partitions & largest upper bound\\
\midrule
$4$ & $23$ & $329$ & $-27813/1000000$\\
$5$ & $216$ & $3{,}757$ & $-14009/1000000$\\
$6$ & $6{,}767$ & $130{,}172$ & $-1/20000$\\
\bottomrule
\end{tabular}
\caption{Certificates. The last column is $\max_I U_I$ over all cells of all entropy profiles.}\label{tab:certificates}
\end{table}

\begin{proof}[Proof of \cref{prop:certificates}]
For each entropy profile of \cref{tab:profiles,tab:six}, bisection as in \cref{sec:partitions} ends with a partition of $[\delta_m,1/2]$ into cells with $U_I<0$; the numbers of cells and the largest $U_I$ are in \cref{tab:certificates}. Each $U_I$ is a rational number obtained by rational arithmetic from the rules of \cref{sec:bernstein,sec:logs,sec:partitions}, so $F_w''\le U_I<0$ on each cell.
\end{proof}

\subsection{Proof of the main theorem}\label{sec:proof}

\begin{proof}[Proof of \cref{thm:main}]
Fix $m\in\{4,5,6\}$ and a positive real vector $w$. For $m=4,5$, \cref{lem:representative} and the listing of integer vectors show that the entropy profile of $w$ is one of the profiles counted in \cref{tab:profiles}. For $m=6$, \cref{prop:sixprofiles} shows that it is one of the profiles counted in \cref{tab:six}. The vector that realizes this profile in the enumeration has the same polynomials $q_s$ in~\eqref{eq:counts}, with multiplicities, as $w$ for all $p$, hence the same entropy function.
By \cref{cor:endpoint}, $F_w''<0$ on $(0,\delta_m]$. By \cref{prop:certificates}, $F_w''<0$ on $[\delta_m,1/2]$. Complement symmetry~\eqref{eq:symmetry} gives $F_w''<0$ on $(0,1)$.

A continuous function on $[0,1]$ with strictly negative second derivative on $(0,1)$ is strictly concave on the closed interval.
\end{proof}

\section{Block results for every number of summands}\label{sec:allm}

\begin{theorem}\label{thm:blocks}
Let $\theta_1,\ldots,\theta_\ell$ be distinct positive real numbers and let $r_1,\ldots,r_\ell$ be positive integers; for each $j$ the variables $B_{j,1},\ldots,B_{j,r_j}$ below form the $j$th \emph{block}. Suppose
\begin{equation}\label{eq:injective}
 \sum_{j=1}^{\ell}\gamma_j\theta_j\ne0
 \quad\text{for every nonzero }\gamma\in\Z^\ell
 \text{ with }|\gamma_j|\le r_j.
\end{equation}
If $B_{j,i}$, $1\le j\le\ell$, $1\le i\le r_j$, are independent Bernoulli$(p)$ variables, then the Shannon entropy of
$\sum_{j=1}^{\ell}\theta_j\sum_{i=1}^{r_j}B_{j,i}$ has strictly negative second derivative on $(0,1)$ and is strictly concave on $[0,1]$. In particular this holds for two positive weight values with irrational ratio, with arbitrary positive multiplicities.
\end{theorem}
\begin{proof}
Put $K_j=\sum_iB_{j,i}$. Every bounded integer vector in \eqref{eq:injective} is the difference of two possible values of $(K_1,\ldots,K_\ell)$. Thus the condition is precisely injectivity of $(K_j)_j\mapsto\sum_j\theta_jK_j$. Entropy is invariant under an injective relabelling, and independence gives
\[
 H\!\left(\sum_j\theta_jK_j\right)
 =H(K_1,\ldots,K_\ell)
 =\sum_jH(\operatorname{Bin}(r_j,p)).
\]
Hillion and Johnson~\cite{HJ} prove the joint concavity conjectured by Shepp and Olkin~\cite{SO}; equal parameters give binomial entropy concavity. The following argument gives a strict second-derivative bound for binomial sizes at least two.

Write $h(p)=-p\log p-(1-p)\log(1-p)$. For a binomial size $\mathfrak n\ge1$ and $K\sim\operatorname{Bin}(\mathfrak n,p)$,
\[
 H(K)=\mathfrak n h(p)-\log(\mathfrak n!)+\mathbb E g(K),
 \qquad g(k)=\log(k!)+\log((\mathfrak n-k)!).
\]
For a function $\phi$ on $\{0,\ldots,\mathfrak n\}$, differentiating the binomial probabilities gives
\[
 \frac{d}{dp}\mathbb E \phi(K)
 =\mathfrak n\mathbb E[\phi(Y+1)-\phi(Y)],\qquad
 Y\sim\operatorname{Bin}(\mathfrak n-1,p).
\]
Indeed, the derivative of the $k$th probability is $\mathfrak n$ times the difference of the $(k-1)$st and $k$th probabilities for $\mathfrak n-1$ summands. Differentiating again gives, for $\mathfrak n\ge2$ and a binomial variable $Z\sim\operatorname{Bin}(\mathfrak n-2,p)$,
\[
 \frac{d^2}{dp^2}\mathbb E g(K)
 =\mathfrak n(\mathfrak n-1)\mathbb E[ g(Z+2)-2g(Z+1)+g(Z)].
\]
The second difference is
\[
 \log\!\left(1+\frac1{Z+1}\right)
 +\log\!\left(1+\frac1{\mathfrak n-Z-1}\right).
\]
Since both denominators are positive, integrating the binomial generating polynomial gives
\[
 \mathbb E\frac1{Z+1}
 =\int_0^1(1-p+p\xi)^{\mathfrak n-2}\,d\xi
 =\frac{1-(1-p)^{\mathfrak n-1}}{(\mathfrak n-1)p}.
\]
Using this identity and its reflection, together with $\log(1+\chi)<\chi$ for $\chi>0$, yields
\begin{align*}
 \frac{d^2}{dp^2}H(\operatorname{Bin}(\mathfrak n,p))
 &<-\frac{\mathfrak n}{p(1-p)}
   +\mathfrak n\frac{1-(1-p)^{\mathfrak n-1}}p
   +\mathfrak n\frac{1-p^{\mathfrak n-1}}{1-p}\\
 &=-\mathfrak n\left(\frac{(1-p)^{\mathfrak n-1}}p+\frac{p^{\mathfrak n-1}}{1-p}\right)<0.
\end{align*}
For $\mathfrak n=1$ the derivative is $h''=-1/[p(1-p)]$. Summing over the blocks shows that the second derivative of the entropy of the sum is negative. Continuity at the endpoints gives strict concavity on the closed interval as in \cref{sec:proof}. For two weight values, an irrational ratio excludes every nonzero integer relation, proving the last assertion.
\end{proof}

\begin{remark}\label{rem:hypotheses}
Concavity in the common parameter can fail for a general Boolean statistic of the Bernoulli variables or for dependent Bernoulli variables with equal marginals. For weighted sums of independent Bernoulli variables, joint concavity in their individual parameters can also fail. The indicator that four independent Bernoulli$(p)$ variables are all equal has success probability $p^4+(1-p)^4$. At $p=1/2$ this probability and its first two derivatives are $1/8$, $0$, $6$, so the second derivative of its entropy is $6\log7>0$.

The heterogeneous example of the MathOverflow user ssss nnnn, given in a comment on Amir's answer and recorded there~\cite{Amir}, has weights $(1,1,2)$ and parameters $(1/5+t,1/5-t,1/2)$. The probabilities at subset sums $0,1,2,3,4$ are
\[
 \left(\frac8{25}-\frac{t^2}2,\frac4{25}+t^2,
 \frac{17}{50}-t^2,\frac4{25}+t^2,\frac1{50}-\frac{t^2}2\right).
\]
Their first derivatives vanish at zero; substitution of their second derivatives $(-1,2,-2,2,-1)$ gives the second derivative $\log(289/256)>0$ of the entropy along this line in parameter space.

For dependent Bernoulli variables with common marginal parameter $p=1/2+t$, take
\[
 (\Pp_{00},\Pp_{01},\Pp_{10},\Pp_{11})
 =\left(\frac14-t+4t^2,\frac14-4t^2,
 \frac14-4t^2,\frac14+t+4t^2\right).
\]
On $|t|\le1/8$ these are probabilities, each at least $3/16$, and both marginals are Bernoulli$(p)$. The sum of the two Bernoulli variables has probabilities $(1/4-t+4t^2,1/2-8t^2,1/4+t+4t^2)$. At zero, \eqref{eq:curvature} gives $16\log2-8>0$.
\end{remark}

\begin{remark}\label{sec:open}

For arbitrary positive real weights and $m\ge7$, the positive-weight concavity question remains open. The endpoint theorems hold for every number of summands, and \cref{thm:blocks} applies when the weighted sum determines how many Bernoulli variables of each weight succeed.

The reduction to integer vectors preserving the signs of differences of two subset sums and the reduction to relation spaces hold for every number of summands. At six summands the coordinate bound $192$ preserving these signs still leaves
$\binom{197}{6}=75{,}176{,}946{,}208$ sorted integer vectors. At seven summands Hadamard's inequality gives $907$, and divisibility of these determinants refines this to $896$, leaving
\[
 \binom{902}{7}=94{,}162{,}890{,}744{,}574{,}080
\]
sorted integer vectors.

The enumeration of relation spaces replaces the listing of bounded integer vectors, but the number of ternary normals grows from $301$ to $966$, and the coordinate-permutation group grows from $720$ to $5{,}040$.

\end{remark}

\section*{Acknowledgments and funding}
The research presented in this paper was supported by the European Research Council (ERC) under the European Union's Horizon Europe research and innovation programme (grant agreement No.~101041711), by the Simons Foundation as part of the Collaboration on the Mathematical and Scientific Foundations of Deep Learning, by Heights Labs, by Convex Nexus Capital, by the Israel Science Foundation (grant number 2258/19), and by the Israel Science Foundation (ISF Grant 4101/25).

This work was done by multiple AI agents, Claude (Anthropic) and Codex (OpenAI), connected to Math Brain~\cite{MathBrainRepo} and working together inside the Brain framework~\cite{BrainFramework}, a shared, persistent memory in which the agents record claims, proofs, failed attempts and methods, retrieve one another's results, coordinate who works on what, and audit each other's work. The method is described in~\cite{MathBrainPaper}.

\section*{Declaration of competing interest}
The author holds equity in Heights Labs and in Convex Nexus Capital and has a management role at Convex Nexus Capital. These entities are included in the funding declaration above.

\section*{Code and data availability}
The \LaTeX\ source, the complete entropy profiles, the rational certificates, and the enumeration and verification programs are available at \url{https://github.com/yspennstate/bernoulli-sum-entropy-concavity}.
The $216$ entropy profiles for $m=5$ are listed with their certificates in the repository. The programs \texttt{certify.py} and \texttt{certify\_six.py} produce the rational partitions and upper bounds. The programs \texttt{replay.py} and \texttt{replay\_six.py} reconstruct the entropy profiles, establish each upper bound and the endpoint constants by rational arithmetic, and establish complete coverage of $[\delta_m,1/2]$. The repository's code README gives the data encodings and the roles of the enumeration and verification programs.

\section*{Declaration of generative AI and AI-assisted technologies}
Codex (OpenAI) proposed the finite reductions and computed the certificates for four, five and six summands; a second Codex run checked every reduction and certificate with separate code. Claude (Anthropic) checked and edited the text. The work rests on the conjecture of Neunh\"auserer~\cite{Neun} and on the formulation and small cases of Amir~\cite{Amir}. It uses the unit-weight theory of Shepp, Olkin, Hillion and Johnson~\cite{SO,HJ}. The coordinate bound adapts the integer-weight method of Muroga, Toda and Takasu~\cite{MTT}, in the account of Servedio~\cite{Servedio}, and its refinement uses the determinant estimates of Ehlich and Wojtas~\cite{Ehlich,Wojtas}. The hyperplane arrangement is that of Chevyrev, Searles and Slinko~\cite{CSS}, and the interval bounds follow the Bernstein-basis account of Farouki~\cite{Farouki}. The author selected the problem, directed the work, and is responsible for the content of the paper and for decisions concerning its dissemination.

\appendix
\section{Notation}\label{app:notation}

\small
\begin{longtable}{p{0.22\textwidth}p{0.43\textwidth}p{0.25\textwidth}}
\toprule
symbol & meaning & defined in\\
\midrule
\endfirsthead
\toprule
symbol & meaning & defined in\\
\midrule
\endhead
\bottomrule
\endfoot
$m$ & number of summands & abstract; \cref{sec:intro}\\
$w=(w_i)$ & positive real weight vector & \cref{sec:intro}\\
$B_i,p$ & independent Bernoulli variables and their common success probability & \cref{sec:intro}\\
$H$ & Shannon entropy, with natural logarithms & \cref{sec:intro}\\
$S_w(p),F_w(p)$ & weighted sum and its entropy & \eqref{eq:sum}\\
$n,\mathfrak a_n,P_n$ & level, vector lengths and law of the recursive overlap construction & \cref{sec:intro}\\
$\mathcal L,\mathcal R$ & placement at the left and right end & \cref{sec:intro}\\
$\omega_i,B_i'$ & weights of the recursive construction and its Bernoulli$(1-p)$ variables & \cref{sec:intro}\\
$s,N_{s,k}$ & subset sum and number of size-$k$ subsets with that sum & \cref{sec:intro}\\
$A,C,k$ & subsets and subset cardinality & \cref{sec:intro}\\
$\Pp$ & probability under the independent Bernoulli law & \cref{sec:setup}\\
$v=(v_i)$ & integer weight vector with a given entropy profile, or with the signs of $w$ & \cref{sec:intro}\\
$\sgn$ & sign, including the value zero & \cref{sec:intro}\\
$a$ & ternary normal vector & \cref{sec:intro}\\
$\delta_4,\delta_5,\delta_6$ & cut points $2^{-12}$, $2^{-14}$, $2^{-15}$ & \cref{sec:intro}\\
$\theta_j,r_j,\ell$ & distinct block weights, multiplicities and number of blocks & \cref{sec:intro}\\
$[m],\mathcal S_w$ & index set $\{1,\ldots,m\}$ and set of subset sums & \cref{sec:setup}\\
$q_s$ & probability of $S_w(p)=s$, a polynomial in $p$ & \eqref{eq:counts}\\
$W$ & total weight $\sum_iw_i$ & \eqref{eq:symmetry}\\
$\R,\Z$ & real numbers and integers & \cref{lem:representative}\\
$M_m$ & integer coordinate bound $\lfloor m^{m/2}\rfloor$ & \cref{lem:representative}\\
$x,\mathcal Q,x^*$ & real vector, feasible polyhedron and a minimizing vertex & proof of \cref{lem:representative}\\
$G,G_i,b,D$ & tight-constraint matrix, column replacement, right side and $\det G$ & proof of \cref{lem:representative}\\
$\mathcal S_v$ & subset sums of the integer weights $v$ & \cref{tab:four}\\
$T_m,L(w),L$ & ternary normals, relation space and a subspace spanned by ternary normals & \cref{sec:six}\\
$\mathbf 1_A$ & indicator vector of a subset & proof of \cref{lem:relations}\\
$\varrho,\mathcal I,P_{\mathcal I}$ & rank, coordinate index set and maximal minors & \cref{sec:six}\\
$J,j_\iota,\iota,\kappa$ & ordered coordinate set, its elements, their index and cofactor vector & \cref{lem:positive}\\
$\mathcal K,E$ & nonnegative kernel cone and support of an extreme ray & proof of \cref{lem:positive}\\
$\eps,\mathbf y$ & small positive number and a vector in $L^\perp$ & proof of \cref{lem:positive}\\
$\mu_p,\mathbf e_j$ & Bernoulli product law and unit point mass at $j$ & proof of \cref{lem:endpoint}\\
$\ell^1$ & sum-of-absolute-masses norm & proof of \cref{lem:endpoint}\\
$w_{\min}$ & smallest weight & proof of \cref{lem:endpoint}\\
$r$ & multiplicity of the smallest weight & \cref{sec:intro}\\
$\eta_m$ & the point $1/(64m^4)$, for every $m\ge3$ & \cref{thm:allendpoint}\\
$k_s,\zeta_s$ & minimum fibre cardinality and leading count & proof of \cref{thm:limit}\\
$f,d,c_k$ & polynomial, degree and Bernstein coefficients & \cref{sec:bernstein}\\
$I,\alpha,\beta$ & rational interval and its endpoints & \cref{sec:bernstein}\\
$\tau,\sigma,c_j^{(\sigma)}$ & split parameter, subdivision level and coefficient array & \cref{sec:bernstein}\\
$\chi,\nu,y,z$ & positive logarithm argument, binary exponent, normalized argument and $(y-1)/(y+1)$ & \cref{sec:logs}\\
$\mathcal N,\mathsf L_{\mathcal N},E_{\mathcal N}$ & series length, logarithm lower bound and remainder bound & \cref{sec:logs}; \eqref{eq:log}\\
$u$ & integration variable in the logarithm series & \cref{sec:logs}\\
$q_-,q_+,\dot q_-,\dot q_+,\ddot q_-,\ddot q_+$ & lower and upper bounds for $q_s,q_s',q_s''$ on a cell & \cref{sec:partitions}\\
$\lambda_-,\lambda_+$ & lower and upper bounds for $\log q_s$ & \cref{sec:partitions}\\
$T_-,V_-,U_I$ & lower product and square bounds, and upper bound for $F_w''$ on $I$ & \cref{sec:partitions}\\
$\widetilde U_I$ & sum of the fibre bounds before the final rounding & \cref{sec:partitions}\\
$\gamma$ & bounded integer relation among block weights & \cref{thm:blocks}\\
$K_j$ & binomial count in block $j$ & proof of \cref{thm:blocks}\\
$\mathfrak n,K,h,g$ & binomial size and count, binary entropy and factorial function & proof of \cref{thm:blocks}\\
$\phi,Y,Z$ & function on the binomial support and auxiliary binomial variables & proof of \cref{thm:blocks}\\
$\mathbb E,\xi$ & expectation and the binomial integration variable & proof of \cref{thm:blocks}\\
$t$ & perturbation parameter in the counterexamples & \cref{rem:hypotheses}\\
\end{longtable}
\normalsize

\begin{thebibliography}{99}

\bibitem{Neun} J.~Neunh\"auserer, \emph{Recursive overlap Bernoulli distributions and an entropy concavity conjecture}, arXiv:2609.05546v2, 2026, \url{https://arxiv.org/abs/2609.05546}.

\bibitem{Amir} Amir, Answer to ``A recursive overlap construction of finite Bernoulli-sum laws and an entropy-concavity conjecture'', \emph{MathOverflow}, question 514729, answer 515134, 11 September 2026, updated 15 September 2026, \url{https://mathoverflow.net/a/515134}.

\bibitem{SO} L.~A.~Shepp and I.~Olkin, Entropy of the sum of independent Bernoulli random variables and of the multinomial distribution, in \emph{Contributions to Probability}, Academic Press, New York, 1981, 201--206, doi:10.1016/B978-0-12-274460-0.50022-9.

\bibitem{HJ} E.~Hillion and O.~Johnson, A proof of the Shepp--Olkin entropy concavity conjecture, \emph{Bernoulli} \textbf{23} (2017), no.~4B, 3638--3649, doi:10.3150/16-BEJ860; arXiv:1503.01570.

\bibitem{Shannon} C.~E.~Shannon, A Mathematical Theory of Communication, \emph{Bell System Technical Journal} \textbf{27} (1948), 379--423, 623--656, doi:10.1002/j.1538-7305.1948.tb01338.x and doi:10.1002/j.1538-7305.1948.tb00917.x.

\bibitem{MTT} S.~Muroga, I.~Toda and S.~Takasu, Theory of majority decision elements, \emph{J. Franklin Inst.} \textbf{271} (1961), no.~5, 376--418, doi:10.1016/0016-0032(61)90702-5.

\bibitem{Servedio} R.~A.~Servedio, Every Linear Threshold Function has a Low-Weight Approximator, \emph{Computational Complexity} \textbf{16} (2007), no.~2, 180--209, doi:10.1007/s00037-007-0228-7.

\bibitem{Matrix} R.~A.~Horn and C.~R.~Johnson, \emph{Matrix Analysis}, second edition, Cambridge University Press, Cambridge, 2012, doi:10.1017/CBO9781139020411.

\bibitem{Farouki} R.~T.~Farouki, The Bernstein polynomial basis: A centennial retrospective, \emph{Computer Aided Geometric Design} \textbf{29} (2012), no.~6, 379--419, doi:10.1016/j.cagd.2012.03.001.

\bibitem{Li} J.~Li, \emph{Multiplicative comparisons of R\'enyi entropies for weighted Bernoulli sums}, arXiv:2609.01529v2, 2026, \url{https://arxiv.org/abs/2609.01529}.

\bibitem{Wang} H.~Wang, \emph{The Sharp R\'enyi and Tsallis Threshold in the Shepp--Olkin Concavity Problem}, arXiv:2609.27433v1, 2026, \url{https://arxiv.org/abs/2609.27433}.

\bibitem{CSS} I.~Chevyrev, D.~Searles and A.~Slinko, On the number of facets of polytopes representing comparative probability orders, \emph{Order} \textbf{30} (2013), no.~3, 749--761, doi:10.1007/s11083-012-9274-0; arXiv:1103.3938.

\bibitem{Ehlich} H.~Ehlich, Determinantenabsch\"atzungen f\"ur bin\"are Matrizen, \emph{Mathematische Zeitschrift} \textbf{83} (1964), 123--132, doi:10.1007/BF01111249.

\bibitem{Wojtas} M.~Wojtas, On Hadamard's inequality for the determinants of order non-divisible by 4, \emph{Colloquium Mathematicum} \textbf{12} (1964), no.~1, 73--83, doi:10.4064/cm-12-1-73-83.

\bibitem{MathBrainRepo} Y.~Shmalo, \emph{Math Brain: public snapshot, code and tests}, 2026, \url{https://github.com/yspennstate/math-brain-public/tree/5427734deee9be643d6dfbdd5e7d4baeeefe4a4a}.

\bibitem{MathBrainPaper} Y.~Shmalo, \emph{Math Brain: shared, auditable memory for multi-agent mathematical research}, October 2026, doi:10.5281/zenodo.24150837, \url{https://github.com/yspennstate/math-brain-public/tree/5427734deee9be643d6dfbdd5e7d4baeeefe4a4a/papers/math-brain-2026-10-24}.

\bibitem{BrainFramework} Y.~Shmalo, \emph{Brain Framework}, public software repository, 2026, \url{https://github.com/yspennstate/brain-framework-public/tree/46b935d26d5a2893f7790b60466387dd8da47ef5}.

\end{thebibliography}
\end{document}
